\cref{thm:connectedcomponent} is the same as Proposition 4 in~\cite{von2007tutorial} (Note the $L$ defined in \cref{thm:connectedcomponent} is equivalent to $L_{sym}$ in~\cite{von2007tutorial}).
We briefly recover the proof below, beginning with a proposition:
\begin{proposition}\cite{von2007tutorial}\label{proposition:expandL}
For every $f\in\mathbb{R}^m$ 
\begin{align}
    f^\top L f &= \frac{1}{2}\sum_{i,j=1}^m w_{i,j}\Big(\frac{f_i}{\sqrt{d_i}} - \frac{f_j}{\sqrt{d_j}}\Big)^2
\end{align}
\end{proposition}
\cref{proposition:expandL} can be verified by simple algebra. 
Now, suppose the graph has $k$ connected components.
We first find $k$ orthonormal eigenvectors.
This can be done by letting $v_1, \ldots, v_k$ be $k$ vectors such that 
\begin{align}
    v_l(j) = \begin{cases} \frac{\sqrt{d_j}}{\sqrt{\sum_{j'\in S_l} d_{j'}}} & j \in S_l \\ 0 & j \not\in S_l \end{cases}
\end{align}
It is easy to verify that $\forall l=1,\ldots,k$, $\|v_l\| = 1$.
For ${l_1}\neq {l_2}$, $\langle v_{l_1}, v_{l_2} \rangle=0$ because $S_{l_2}$ and $S_{l_1}$ are disjoint.
Finally, the $i$-th entry of $\sqrt{\sum_{j'\in S_l} d_{j'}} Lv_l$ is 0 if $i\not\in S_l$, and 
\begin{align}
    \sqrt{d_i} - \frac{1}{\sqrt{d_i}} \sum_{j}w_{ij} \frac{1}{\sqrt{d_i}}\sqrt{d_i} = 0
\end{align}
if $i\in S_l$ (because $L = I - D^{-\frac{1}{2}} W D^{-\frac{1}{2}}$).

Now, we argue that we cannot find another vector $v$ that is a zero eigenvector.
By \cref{proposition:expandL} $v$ must be be a scaled version of $v_l$ on component $S_l$.
w.l.o.g., assume $v(j)\neq 0$, with $j\in S_l$.
Then, by \cref{proposition:expandL}, $\exists c$ such that $\forall {j'}\in S_l, v(j') = c v_l(j')$.
This means $\langle v, v_l\rangle \neq 0$.
Thus, we cannot find another zero eigenvector. 
\qed


% Reference https://timroughgarden.org/s17/l/l11.pdf Section 2.1